\section{Practical recommendations}
\label{sec:recommendations}

The controlled results above translate into concrete guidance. We separate
recommendations for practitioners choosing or applying a tool today from
recommendations for those designing the next generation of models.

\subsection{For practitioners}

\begin{table*}[t]
\centering
\caption{\textbf{Use-case-specific recommendations.} ``long-$k$ model'' denotes
an overlapping long-$k$-mer neural classifier; ``6-mer GFM'' denotes a
pre-trained NT-v2-style model with non-overlapping 6-mer tokenization.}
\label{tab:reco}
\begin{tabularx}{\textwidth}{@{}p{0.30\textwidth}X@{}}
\toprule
Use case & Recommendation \\
\midrule
Read-level genus classification &
Prefer an overlapping long-$k$ model (87\% vs 67\%); a 6-mer GFM is a distant
second. \\[2pt]
Sample-level abundance (composition) &
A 67\%-accurate model is often sufficient ($r=0.99$) and can beat Kraken\,2 by
avoiding abstention bias---but validate for systematic bias on your community. \\[2pt]
Presence/absence detection, esp.\ low-abundance taxa &
Use Kraken\,2 (or another exact-$k$-mer tool); current 6-mer neural detectors
are weak. Apply a calibrated abundance threshold above the $\sim$0.3\% noise
floor. \\[2pt]
Species-level classification &
6-mer GFMs are insufficient; use a long-$k$ tokenizer or a hierarchical router,
or an exact-$k$-mer tool when the target is in-database. \\[2pt]
Handling class imbalance &
Do \emph{not} aggressively genus-balance if you report abundance: it lifts macro
metrics but harms per-read accuracy and abundance fidelity ($r\!:0.99\!\to\!0.86$). \\[2pt]
Choosing a GFM &
Do not select on backbone size alone; the tokenizer decides whether taxonomic
signal survives. Ask what $k$ and what stride the model uses. \\
\botrule
\end{tabularx}
\end{table*}

\subsection{For model builders}

The single highest-leverage design choice is tokenization. Future microbial
foundation models intended for short-read taxonomy should be \emph{pre-trained}
with overlapping long-$k$ tokenization (e.g.\ 12--13-mers), combining the
discriminative power that overlapping long $k$-mers provide with the
representational quality that pre-training adds---the two independent gains we
quantified ($+20$ pp from tokenization, $+13$ pp from pre-training). Because a
naive long-$k$ vocabulary is enormous (a 13-mer vocabulary implies a
multi-billion-parameter embedding table; the ``$\sim$5M-parameter''
MetaTransformer figure counts only the encoder and understates true cost),
practical realizations will need vocabulary factorization, hashing, or hybrid
tokenization---an open engineering problem we regard as the key next step.
Pre-training corpora should draw on the hundreds of thousands of microbial
genomes now catalogued \citep{Parks2022,Almeida2019} rather than
eukaryote-dominated data.
